% Checked supporting fragment. This replaces any identical-map ordering
% assertion for independently addressed PG and PA exponential tables.
% It is a project-specific derivation, not a claim of theoretical priority.
\section{Independent payoff-table precision and projected encoding}
\label{sec:raw-asymmetric-projection}

I allow the two exponential look-up tables in the common-normal payoff
difference to use different address widths. The two represented maps can
be individually monotone without preserving their mutual order. The
financial target remains the original no-arbitrage conditional correction;
the following projection makes the implemented asymmetric difference a
valid probability without assuming an ordering that has not been proved.

Throughout this section, $K=100$, $H=1$, $\lambda>0$, $q=28$ and
$s=2^q$. Define
\begin{equation}
 f(x)=\min\{1,100(e^x-1)_+\},\qquad
 \Pi(t)=\min\{1,\max\{0,t\}\},\qquad
 \ell=\log(101/100).
\end{equation}
Since $\ell<1/64$, the exact bounded-argument decomposition is
\begin{equation}
 f(x)=\min\{1,100[\exp(\min\{1/64,x_+\})-1]_+\}.
\end{equation}
Indeed, the two expressions agree for $x\leq0$, for
$0<x<\ell$, and for $x\geq\ell$. The last case follows because
the clamped argument is at least $\ell$ whenever the unclamped
argument is at least $\ell$. Consequently, the exponential only needs
its certified domain $[0,1/64]$, with
$1\leq e^t\leq64/63$ there, the last inequality following from
$e^t\leq(1-t)^{-1}$ for $0\leq t<1$.

\begin{lemma}[Exact projected integer difference]
\label{lem:raw-asymmetric-projection}
Suppose every unsigned operand and intermediate product satisfies the
retained domain and overflow certificate. Let $a,b$ be the represented
payoff words, decoded in $s^{-1}\mathbb Z\cap[0,1]$. Then the emitted
operation \texttt{clamp(positive\_sub($a$,$b$),1)} decodes exactly to
$\Pi(a-b)$. In particular, its range is $[0,1]$ even if $a<b$.
The uniform $q$-bit comparator encodes this value as a flag probability
exactly.
\end{lemma}
\begin{proof}
Let the word width be $w$ and write $A=sa$, $B=sb$. The subtraction
primitive produces $A-B$ modulo $2^w$ in temporary workspace. The
comparison flag is one exactly when $B<A$. Controlled copies of the
temporary difference to a zero output word therefore return $A-B$ when
$A>B$, and zero when $A\leq B$. When $A>B$, the retained unsigned
bounds ensure $0<A-B<2^w$, so the copied word is the ordinary integer
difference. When $A<B$, the modular subtraction is never copied and is
uncomputed with the comparison workspace. The temporary modular word
does not become a financial value. The subsequent clamp at the exactly
representable integer $s$ gives $\min\{s,\max\{0,A-B\}\}$.
The guarded operands are nonnegative and fit their registers; the
retained pre-mask magnitude guard also excludes an ambiguous overflow.
Thus decoding gives exactly $\Pi(a-b)$. For an output integer
$C\in\{0,\ldots,s\}$, exactly $C$ labels in the uniform set
$\{0,\ldots,s-1\}$ satisfy $U<C$. The good-flag probability is
$C/s$, including both endpoints. The clean inverse restores the
workspace after the flag is copied.
\end{proof}

The projection is nondecreasing and $1$-Lipschitz. For example, if
$u\leq v$, then its piecewise-linear definition gives
$0\leq\Pi(v)-\Pi(u)\leq v-u$. Reversing the arguments gives the
absolute-value inequality. For every $\Delta\geq0$,
$h_{m,\Delta}(n)=f(m+\Delta+\lambda n)-f(m+\lambda n)$ lies in
$[0,1]$, so $\Pi(h_{m,\Delta}(n))=h_{m,\Delta}(n)$.

\begin{lemma}[Separate scalar errors survive projection]
\label{lem:raw-asymmetric-scalar}
Let $\widehat E_G$ and $\widehat E_A$ be the two guarded dyadic
exponential implementations, with independently certified bounds
\begin{equation}
 \sup_{t\in[0,1/64]\cap s^{-1}\mathbb Z}
 |\widehat E_i(t)-e^t|\leq e_{\exp,i},\qquad i\in\{G,A\}.
\end{equation}
Let $\widehat f_i$ denote the exact emitted clamps and positive
subtractions around $\widehat E_i$. For any guarded decoded arguments
$x_G,x_A$,
\begin{equation}
 \left|\Pi(\widehat f_A(x_A)-\widehat f_G(x_G))
          -\Pi(f(x_A)-f(x_G))\right|
 \leq100(e_{\exp,G}+e_{\exp,A}).
\end{equation}
\end{lemma}
\begin{proof}
The constants zero, one and $1/64$ are exactly representable at
$q=28$. All payoff clamps and positive subtractions implement their
decoded real maps exactly, conditional on the whole-graph guards.
Integer scaling by $100$ has no fixed-point floor. Since these clamps
are $1$-Lipschitz, $|\widehat f_i(x)-f(x)|\leq100e_{\exp,i}$.
The triangle inequality for the difference and the $1$-Lipschitz
property of $\Pi$ give the claim. No comparison between
$\widehat f_G$ and $\widehat f_A$ is required. Each $e_{\exp,i}$
is the scalar's \texttt{output\_error\_upper}, comprising nodal
rounding, interpolation error and interpolation floor. Its propagated
argument-error charge is excluded here because that charge is treated
separately below.
\end{proof}

Fix a realised outer-grid coordinate $y$. Let $m=m_G(y)$,
$\Delta=\log T(y)\geq0$, and let the inherited shifted node have
ideal value $M=m+B_G$ and represented value $\widehat M$.
Write $\widehat\Delta$ for the represented logarithm. The guarded
graph gives
\begin{equation}
 |\widehat M-M|\leq e_m,\qquad
 |\widehat\Delta-\Delta|\leq e_\Delta,\qquad
 \widehat\Delta\geq0.
\end{equation}
These errors treat the realised outer coordinates as exact inputs.
Outer Gaussian-node rounding is already charged by the inherited
continuous-to-finite financial transfer and is not charged again here.
Let $C=J-B_G$, $O=J+\lambda L$ and let their represented constants
be $\widehat C$, $\widehat O$. If
$|\widehat C-C|\leq e_C$ and $|\widehat O-O|\leq e_O$, set
\begin{equation}
 \widehat m=\widehat M-B_G+(\widehat C-C)-(\widehat O-O),
 \qquad E_m=e_m+e_C+e_O.
\end{equation}
Thus $|\widehat m-m|\leq E_m$.

Let $(I_\lambda,d_\lambda)$ be the actual integer coefficient and
drop count of the scaling emitter, and put
\begin{equation}
 \widehat\lambda=I_\lambda2^{-d_\lambda},\qquad
 E_s=2L|\widehat\lambda-\lambda|
       +s^{-1}\mathbf1_{\{d_\lambda>0\}}.
\end{equation}
For a realised quantile $\widetilde N\in[-L,L]$, the exact decoded
scale node is
$\widehat\lambda(\widetilde N+L)+\rho(\widetilde N+L)$,
where $|\rho|\leq s^{-1}$ when $d_\lambda>0$, and $\rho=0$
otherwise. The guarded endpoint additions are exact integer additions.
Consequently their conceptual signed arguments are
\begin{align}
 \widehat X_G&=\widehat m+\lambda\widetilde N+
                    \zeta(\widetilde N),\qquad |\zeta|\leq E_s,\label{eq:raw-asymmetric-arguments}\\
 \widehat X_A&=\widehat X_G+\widehat\Delta.\nonumber
\end{align}
These are mathematical signed differences used to identify the financial
target. The circuit obtains their positive parts by the guarded unsigned
construction of Lemma~\ref{lem:raw-asymmetric-projection}; it does not
store an unguarded signed log argument.

\begin{theorem}[Conditional error for independently addressed tables]
\label{thm:raw-asymmetric-conditional}
Assume the whole-graph certificate passes and the realised finite normal
law has Kolmogorov discrepancy at most $D_N$ from the standard normal
law. Let $\nu_0\geq\sup_x\phi(x)$ and let $p(y)$ be the implemented
good-flag probability conditional on the realised outer coordinate. Then
\begin{equation}
 |p(y)-R_1(y)|\leq\eta_{\rm smooth}:=
 \min\left\{1,100(e_{\exp,G}+e_{\exp,A})+202E_s+2D_N+
       \frac{\nu_0}{\lambda}(2E_m+e_\Delta)\right\}.
\end{equation}
The minimum of this bound and an independently valid pointwise
graph/CDF bound is a valid uniform conditional certificate.
\end{theorem}
\begin{proof}
The payoff is absolutely continuous, with almost-everywhere derivative
$f'(x)=100e^x\mathbf1_{(0,\ell)}(x)$, and is globally
$101$-Lipschitz. Replacing the two arguments in
\eqref{eq:raw-asymmetric-arguments} by
$\widehat m+\lambda\widetilde N$ and
$\widehat m+\widehat\Delta+\lambda\widetilde N$ changes the exact
payoff difference by at most $202E_s$ per normal label. Projection
cannot enlarge this error and fixes the latter difference because
$\widehat\Delta\geq0$. Lemma~\ref{lem:raw-asymmetric-scalar}
adds at most $100(e_{\exp,G}+e_{\exp,A})$.

The remaining finite-inner-law target is
$\mathbb E h_{\widehat m,\widehat\Delta}(\widetilde N)$. The
parameters $\widehat m$ and $\widehat\Delta$ depend only on outer
labels, not on the inner normal label. Each summand of $h$ is
continuous and nondecreasing with total variation one, so
$\operatorname{TV}(h)\leq2$. Stieltjes integration by parts gives
\begin{equation}
 \mathbb E h(\widetilde N)-\mathbb E h(N)
   =-\int_{\mathbb R}(F_{\widetilde N}(x)-\Phi(x))\,dh(x).
\end{equation}
The boundary term vanishes because $h$ is bounded and both CDFs tend
to the same limits at infinity. The continuous function $h$ assigns
no atom to its derivative measure, so the identity accommodates
atoms of the finite law and duplicate quantiles. Its absolute value
is at most $2D_N$. Normal rounding is charged solely within $D_N$.

Define $P_\lambda(t)=\mathbb E f(t+\lambda N)$. The nonnegative
measure $df$ has mass
$100(e^\ell-1)=1$. Gaussian convolution gives
\begin{equation}
 P_\lambda'(t)=\int_{\mathbb R}
       \frac1\lambda\phi\left(\frac{x-t}{\lambda}\right)df(x),
 \qquad 0\leq P_\lambda'(t)\leq\frac{\nu_0}{\lambda}.
\end{equation}
Differentiation is justified by the bounded smooth Gaussian kernel
against the finite measure $df$; equivalently, one may use the
bounded derivative of $f$ and dominated convergence. Thus the two
continuous endpoint perturbations contribute at most
$\nu_0(2E_m+e_\Delta)/\lambda$. Adding the independently justified
terms proves the uncapped bound. Both $p(y)$ and $R_1(y)$ lie in
$[0,1]$, so capping the error at one is valid. The minimum of two
valid uniform upper bounds remains valid.
\end{proof}

\begin{proposition}[Prospective common-shift refinement]
\label{prop:raw-common-shift-refinement}
For $R(m,\Delta)=P_\lambda(m+\Delta)-P_\lambda(m)$ with
$\Delta\geq0$, the continuous outer perturbation above also satisfies
\begin{equation}
 |R(\widehat m,\widehat\Delta)-R(m,\Delta)|
 \leq\frac{\nu_0}{\lambda}(E_m+e_\Delta).
\end{equation}
If $\nu_1\geq\sup_x|\phi'(x)|$ and a certified uniform bound
$\Delta\leq\Delta_*$ is available, the common-shift coefficient
can be replaced by
\begin{equation}
 \min\left\{\frac{\nu_0}{\lambda},
              \frac{\nu_1\Delta_*}{\lambda^2}\right\}E_m
       +\frac{\nu_0}{\lambda}e_\Delta.
\end{equation}
\end{proposition}
\begin{proof}
The previous convolution identity bounds $P_\lambda'$ in
$[0,\nu_0/\lambda]$. Hence
$|\partial_mR|=|P_\lambda'(m+\Delta)-P_\lambda'(m)|
\leq\nu_0/\lambda$, whereas
$|\partial_\Delta R|=P_\lambda'(m+\Delta)\leq\nu_0/\lambda$.
Split the perturbation into a gap change at $\widehat m$ followed by
a common-shift change at the true nonnegative gap $\Delta$, and
integrate these derivative bounds along the two intervals. Furthermore,
differentiating the Gaussian convolution against the probability
measure $df$ gives $|P_\lambda''|\leq\nu_1/\lambda^2$.
Integrating $P_\lambda''$ between $m$ and $m+\Delta$ bounds
$|\partial_mR|$ also by $\nu_1\Delta/\lambda^2$.
The uniform upper bound $\Delta_*$ gives the refined coefficient.
\end{proof}

This refinement is an analytic error bound, not an automatic improvement
in the full estimation plan. It must be selected and recorded in a
protocol before any resource outcomes using it are generated. The
financial coefficient, external outer-law mass, continuous transfer,
mean prior, estimator accuracy, confidence and full clean circuit costs
must be retained or rebuilt consistently. The projection argument
justifies asymmetric scalar precision without claiming cancellation of
the two distinct table errors, a new amplitude-estimation query law, or
a physical quantum advantage.
